\documentclass[../main]{subfiles}
\begin{document}

\section{Teorema de Continuidad en Hidrodinámica}

El teorema de continuidad en hidrodinámica es un principio fundamental que establece la conservación de la masa en un flujo de fluido. Este principio es esencial para comprender y analizar diversos fenómenos en fluidos, desde el flujo en tuberías hasta la aerodinámica de aviones.

\subsection{Fundamentos Teóricos}

El teorema de continuidad establece que, en un flujo estacionario e incompresible de un fluido, la masa que entra en una sección transversal del flujo debe ser igual a la masa que sale de esa sección transversal. Matemáticamente, esto se expresa como:

\info{Teorema de Continuidad}{
En un flujo estacionario e incompresible, la ecuación de continuidad es:
\begin{equation}
A_1 \cdot v_1 = A_2 \cdot v_2
\end{equation}
donde:
\begin{itemize}
    \item $A_1$ y $A_2$ son las áreas de las secciones transversales en los puntos 1 y 2 del flujo, respectivamente.
    \item $v_1$ y $v_2$ son las velocidades del flujo en los puntos 1 y 2, respectivamente.
\end{itemize}
}

El teorema de continuidad se basa en el principio de conservación de la masa, que establece que la masa total de un sistema cerrado permanece constante con el tiempo, siempre que no haya flujo de masa hacia adentro o hacia afuera del sistema.

\subsection{Aplicaciones}

El teorema de continuidad tiene una amplia gama de aplicaciones en la ingeniería y la física, incluyendo:

\begin{itemize}
    \item \textbf{Flujo de fluidos en tuberías}: Se utiliza para calcular el caudal volumétrico de un fluido en una tubería, que es fundamental para el diseño y la operación de sistemas de distribución de agua, sistemas de ventilación, sistemas de climatización, entre otros.
    \item \textbf{Aerodinámica}: Se aplica para analizar el flujo de aire alrededor de cuerpos sólidos, como alas de aviones, para determinar la generación de sustentación y resistencia.
    \item \textbf{Hidráulica}: Se emplea para estudiar el flujo de agua en canales abiertos, ríos y embalses, para el diseño de estructuras hidráulicas como presas, compuertas y sistemas de riego.
    \item \textbf{Meteorología}: Se utiliza para modelar y predecir el movimiento de masas de aire en la atmósfera, así como el comportamiento de los océanos y los ríos.
\end{itemize}

\subsection{Preguntas para Investigar}

\begin{enumerate}
    \item ¿Cuál es la diferencia entre un flujo estacionario y un flujo no estacionario?
    \item ¿Cómo se puede aplicar el teorema de continuidad en la industria alimentaria?
    \item ¿Qué ocurre si el flujo es compresible en lugar de incompresible?
    \item ¿Cuál es la importancia del teorema de continuidad en la medicina, especialmente en el estudio de la circulación sanguínea?
    \item ¿Cómo se puede utilizar el teorema de continuidad para determinar la velocidad del agua en un río?
    \item ¿Cuáles son las limitaciones del teorema de continuidad en la modelización de sistemas complejos de fluidos?
    \item ¿Cómo se puede aplicar el teorema de continuidad en la industria aeroespacial?
    \item ¿Qué tipo de instrumentos se utilizan para medir el flujo de fluidos en aplicaciones industriales?
    \item ¿Cuáles son las implicaciones del teorema de continuidad en el diseño de sistemas de alcantarillado urbano?
    \item ¿Cómo se relaciona el teorema de continuidad con la ley de conservación de la energía en fluidos?
\end{enumerate}
\subsection{Problemas Prácticos}

\begin{enumerate}
    \item Una tubería de 10 cm de diámetro tiene un flujo de agua con una velocidad de 2 m/s. ¿Cuál es el diámetro de una tubería en la que el flujo de agua tiene una velocidad de 1 m/s?
    \item Si un ventilador mueve 100 m\textsuperscript{3} de aire por minuto a través de una sección transversal de 2 m\textsuperscript{2}, ¿cuál es la velocidad del aire en m/s?
    \item Un río tiene una sección transversal de 50 m\textsuperscript{2} y una velocidad de flujo de 2 m/s. Si se reduce la sección transversal a 25 m\textsuperscript{2}, ¿cuál será la velocidad del flujo?
    \item Se necesita suministrar 500 litros de agua por minuto a través de una tubería de 20 cm de diámetro. ¿Cuál debe ser la velocidad del flujo en m/s?
    \item Una aeronave vuela a través de una corriente de aire a una velocidad de 200 m/s. Si el área frontal de la aeronave es de 100 m\textsuperscript{2}, ¿cuál es el volumen de aire que pasa por la aeronave por segundo?
    \item Un conducto de ventilación tiene una sección transversal de 1 m\textsuperscript{2} y una velocidad de aire de 5 m/s. Si se necesita aumentar la velocidad del aire a 10 m/s, ¿cuál debe ser la nueva sección transversal?
    \item Un sistema de riego necesita suministrar 1000 litros de agua por hora a través de una tubería de 15 cm de diámetro. ¿Cuál debe ser la velocidad del flujo en m/s?
    \item Un río fluye a una velocidad de 3 m/s y tiene una sección transversal de 100 m\textsuperscript{2}. Si se ensancha el río y se duplica su sección transversal, ¿cuál será la nueva velocidad del flujo?
    \item Un tanque tiene una sección transversal de 5 m\textsuperscript{2} y está lleno de agua hasta una altura de 3 metros. Si se abre una válvula en la parte inferior del tanque, ¿cuál será la velocidad del flujo de agua?
    \item Una bomba de agua suministra 200 litros de agua por segundo a través de una tubería de 30 cm de diámetro. ¿Cuál es la velocidad del flujo en m/s?
\end{enumerate}

\end{document}
